\documentclass[conference]{IEEEtran}
\usepackage[utf8]{inputenc}
\usepackage[T1]{fontenc}
\usepackage{cite}
\usepackage{amsmath,amssymb}
\usepackage{graphicx}
\usepackage{booktabs}
\usepackage{url}

\title{Measuring Validator Utility in Generate--Verify--Repair NetOps Pipelines: a Confirmatory Paired Study with Shared Initial Candidates}

\author{
\IEEEauthorblockN{Autonomous AI Researcher}
\IEEEauthorblockA{CNSM 2026 Agentic AI Researcher Track}
}

\begin{document}

\maketitle

\begin{abstract}
We preregistered and executed a paired confirmatory study (AC-2026-03) to evaluate whether a multidimensional validator metric suite predicts end-to-end agentic NetOps success better than a single verifier pass rate. Using a synthetic intent\ensuremath{\rightarrow}configuration generator and a hosted generate--validate--repair (GVR) adapter under shared\_initial\_candidate semantics, we ran 200 paired tasks (one shared initial generation per task reused for baseline and guarded). The deterministic validator (deterministic\_netops\_validator\_v5, v5.0) judged every sampled initial candidate valid; no repairs were applied. Baseline = 200/200, guarded = 200/200, paired difference = 0.00 (bootstrap 95\% CI [0.00, 0.00], exact McNemar two-sided p = 1.0; bootstrap\_seed = 1729, resamples = 10000). The observed ceiling invalidated the preregistered power assumptions (targeted 10 percentage-point paired improvement) and rendered the primary confirmatory test uninformative for the originally planned effect. We report preregistered design choices, precise execution accounting, representative validator-trace excerpts, quantitative analysis, failure-mode diagnostics, and artifact-grounded prescriptions for follow-up preregistered experiments (fault injection, multi-sample generation, multi-validator comparison) required to estimate validator coverage and repair value.
\end{abstract}

\section{Introduction}
Automating NetOps workflows with generative models requires reliable mechanisms to avoid harmful, silent, or wrong-but-plausible actions. A common operational pattern is generate--verify--repair (GVR): a generator (LLM) proposes a candidate configuration or plan, a deterministic validator mechanically checks correctness and policy constraints, and, when necessary, a repair module synthesizes corrected candidates or triggers human review. Prior demonstrations of GVR-like architectures in code and systems motivate their adoption in NetOps, but systematic, preregistered measurements of validator coverage (false-positive/false-negative rates), detection latency, and predictive usefulness for end-to-end guarded success are scarce [1].

We preregistered study AC-2026-03 with a paired design that isolates the effect of downstream validator-triggered repair under shared\_initial\_candidate semantics: for each task the generator produced a single initial candidate that was reused by both the baseline (no repair) and guarded (validator+repair permitted) conditions. Under these semantics any paired difference can only arise downstream from validator-triggered repair, not from independent generator sampling. The primary estimand was the paired\_success\_rate\_difference\_guarded\_minus\_baseline and the preregistered confirmatory hypothesis (H1) expected that a multidimensional validator-metric suite would explain more variance in end-to-end success than a single verifier pass rate. Our main contribution is a fully preregistered, artifact-archived experiment and an exact, transparent report of a degenerate confirmatory outcome (ceiling) with concrete, artifact-grounded prescriptions for follow-up designs that can measure validator utility.

\section{Related Work}
Work across code, systems, and NetOps communities has proposed and prototyped generate--verify--repair or closed-loop verifier architectures to reduce stochastic generator error and enable deterministic acceptance criteria [1]. Recent NetOps and AIOps benchmark efforts document substantive operational failure modes (ticket noise, false premises, and wrong-but-plausible outputs) that complicate end-to-end automation and motivate insertion of verifier layers [2],[3]. Independent system reports and preprints propose self-healing orchestrators, auditable decision traces, and verifier-guided repair to reduce silent or action-triggering failures in tool-augmented LLM systems [4]. Surveys of LLMs for telecommunications and NetOps discuss adaptation and verification trade-offs and motivate metricized evaluation of verification cost versus nominal correctness [5].

These verified records motivate a measurement focus on validator soundness and operational impact but do not provide a standardized, empirically estimable validator-metric suite tailored to NetOps GVR flows. The present preregistered experiment was designed to begin filling that measurement gap under strict preregistration and archived-execution practices; when interpreted against the cited literature we treat prior demonstrations as architectural and motivational rather than as establishing empirical coverage for NetOps validators.

\section{Methodology}
Preregistration and primary design. The study preregistration (AC-2026-03) specified a paired confirmatory design using the hosted\_netops\_gvr\_v1 adapter family with generation\_semantics = shared\_initial\_candidate and initial\_generation\_calls\_per\_task = 1. Planned task\_count = 200 pairs (planned\_episode\_count = 400), planned maximum\_model\_calls = 400, and model\_scope = gpt-5-mini via the provider bundle recorded as 'openai\_responses'. The preregistered primary estimand was paired\_success\_rate\_difference\_guarded\_minus\_baseline; the primary test was an exact two-sided McNemar with paired bootstrap percentile confidence intervals (bootstrap\_resamples = 10000, bootstrap\_seed = 1729). Targeted detectable effect was an absolute 0.10 paired improvement with \ensuremath{\alpha} = 0.05 and \textasciitilde{}80\% power under conservative planning assumptions.

Generator, tasks, and constraints. Tasks were drawn deterministically from deterministic\_netops\_task\_generator\_v5 (task\_family = intent\_configuration\_repair\_v1). Each task encodes: an initial\_state, an expected\_final\_state, an explicit intent string, allowed\_fields/allowed\_touched\_fields, workflow\_policies (e.g., minimum\_active\_in\_groups, must\_be\_down\_for\_fields), and a required\_command\_sequence. Task selection followed a preregistered deterministic index rule (challenging\_workflow\_stress\_v1). Contamination/exclusion rules (duplicate SHA-256 and embedding-similarity > 0.95) were preregistered; the execution/analysis artifacts record that no contamination exclusions occurred (analysis/contamination\_summary.csv empty).

Model and sampling configuration. The declared model version in execution/model\_configuration.json is "gpt-5-mini"; execution accounting records model\_calls\_used = 200 (one shared initial generation per pair) and maximum\_model\_calls = 400 as preregistered. Important artifact-grounded note: execution/model\_configuration.json records temperature = null (temperature = null/unset). Null/unset is explicitly not evidence of deterministic provider-side sampling and does not imply temperature = 0; provider sampling variability remains possible and is recorded in provider-call audit fields. Provider-call audit in deterministic\_reconciliation reports hosted\_model\_call\_event\_count = 200 and stage\_counts.shared\_initial\_generation = 200, consistent with the preregistered single shared generation per task.

Validator and scoring. The deterministic validator used was deterministic\_netops\_validator\_v5 (validator\_version = 5.0). The validator implements mechanistic intent-constraint checks and emits per-episode fields: normalized\_configuration, parsed\_command\_count, required\_command\_count, observed\_touched\_field\_count, violation\_count/violation\_codes, valid (boolean), repair\_applied (boolean), and repair artifacts when applicable. The preregistered binary scoring rule (full\_intent\_constraint\_validation) uses validator.valid as the success label for the confirmatory estimand. The pipeline permitted at most one repair call when validator rejected a candidate; repair events would be recorded with repair\_applied = true and repair\_provider\_trace fields.

Execution, exclusions, and missingness. The preregistered maximum\_attempts\_per\_call = 1 (no manual retries). The execution manifest reports completed\_episode\_count = 400, failed\_episode\_count = 0, and complete\_pair\_count = 200. The preregistered missingness and contamination plans were followed: confirmatory analyses used complete pairs only; artifacts include analysis/missingness\_summary.csv and analysis/contamination\_summary.csv. All raw responses, per-episode validator traces, provider-call traces, and analysis artifacts are archived in the execution bundle and referenced in this paper.

\section{Results}
Execution accounting and primary outcomes. The adapter completed the run exactly as preregistered: planned\_episode\_count = 400, completed\_episode\_count = 400, failed\_episode\_count = 0, and model\_calls\_used = 200 (one shared initial generation per pair). Provider-call audit recorded hosted\_model\_call\_event\_count = 200, completed\_hosted\_model\_call\_event\_count = 200, cache\_hit\_count = 0, cache\_miss\_count = 200, and stage\_counts.shared\_initial\_generation = 200 (deterministic\_reconciliation.json). Contamination and missingness artifacts confirm no exclusions or missing episodes (analysis/contamination\_summary.csv empty; analysis/missingness\_summary.csv reports complete\_pairs = 200 and zeros for all missingness categories).

Primary confirmatory result and exact accounting. The deterministic validator judged every sampled shared initial candidate as valid; no repairs were applied in any episode. Analysis artifacts report baseline\_success\_count = 200/200 (success\_rate = 1.0) and guarded\_success\_count = 200/200 (success\_rate = 1.0). The paired contingency table (analysis/paired\_contingency\_table.csv) is n\_11 = 200, n\_10 = 0, n\_01 = 0, n\_00 = 0. Consequently the paired\_success\_rate\_difference\_guarded\_minus\_baseline = 0.00. The prespecified exact McNemar two-sided test yields p-value = 1.0. The paired bootstrap percentile 95\% confidence interval (bootstrap\_resamples = 10000, bootstrap\_seed = 1729; analysis/analysis\_log.jsonl) is [0.00, 0.00]. These numeric artifacts (analysis/condition\_summary.csv, analysis/paired\_contingency\_table.csv, analysis/analysis\_log.jsonl) are archived and match the reported values.

Representative per-episode validator diagnostics (artifact-grounded). To illustrate why no repairs were triggered, Table and scoring artifacts contain the validator trace fields for every episode; a compact example (execution/scoring/task-000004-guarded.json) shows the key mechanistic outputs that were recorded for each episode:

pair\_id: pair-000004
  normalized\_configuration: "interface edge2 admin down\textbackslash{}ninterface edge2 vlan 40\textbackslash{}ninterface edge2 admin up\textbackslash{}ninterface edge1 admin down"
  parsed\_command\_count: 4
  required\_command\_count: 4
  observed\_touched\_field\_count: 3
  valid: true
  repair\_applied: false
  validator\_id: deterministic\_netops\_validator\_v5
  validator\_version: 5.0

Equivalent validator\_trace.validation\_after.valid = true and repair\_applied = false hold for every archived episode; the repair\_provider\_trace fields are null for all episodes (execution/scoring/*.json entries list repair\_provider\_trace\_path = null). Representative tasks in the artifact\_examples (task-000004, task-000005, task-000006) likewise show parsed\_command\_count equal to required\_command\_count and valid = true, demonstrating that, for these sampled seeds, generator outputs satisfied the validator's mechanistic constraints exactly.

Secondary and exploratory diagnostics (uninformative under ceiling). Because repair\_applied = false for all episodes, secondary metrics that depend on rejection or repair events are degenerate in this execution: per-validator false-positive or false-negative rate estimates cannot be computed from observed rejections; detection-latency and repair-invocation distributions are empty by design. The archived raw\_results.jsonl and per-episode validator traces nevertheless provide the full data required to compute these metrics under alternative scoring rules or after applying the follow-up prescriptions (fault injection or multi-sample generation) described elsewhere in this paper.

Artifact-grounded diagnostic interpretation. Under shared\_initial\_candidate semantics any paired difference could only arise from validator-triggered repair. The degenerate all-valid outcome observed here is explicable from artifact-grounded evidence in three linked facts: (1) the deterministic task generator and the recorded generation seeds produced initial candidates that satisfied the validator's mechanistic intent-constraint checks (as shown by per-episode validation\_after summaries); (2) the single-sample-per-task sampling regime removed generator variability as a source of disagreement by design (execution/model\_configuration.json and deterministic\_reconciliation.json record one shared generation per task); (3) mechanistic validators operate on structural and policy checks and can be blind to semantic intent misalignment (correct-by-structure yet incorrect-by-intent) absent a separate semantic-intent model or operator feedback. These jointly explain why no repairs occurred and why the confirmatory test could not detect the preregistered 10-percentage-point improvement target.

Reproducibility pointers and archived artifacts. All numeric claims above are reproducible from the archived artifacts: execution/raw\_results.jsonl (input\_results\_sha256 = 23f4cd01\allowbreak{}0f3a7419\allowbreak{}eb21ef9b\allowbreak{}8caf4dbe\allowbreak{}7730552e\allowbreak{}aa3f3eaa\allowbreak{}22d2068c\allowbreak{}f3b5d4e3), execution/model\_configuration.json (sha256 = a40e96e2...), analysis/paired\_contingency\_table.csv (sha256 = 7bc1bd2a...), analysis/condition\_summary.csv (sha256 = 1cb072b4...), analysis/analysis\_log.jsonl (sha256 = dd2dbf18...), and representative per-episode scoring files (execution/scoring/task-*.json). These files show the per-episode validator.valid flags, repair\_applied markers, shared\_initial\_candidate content and SHA-256s, and provider-call audit entries used to produce the summary statistics above. Because the observed ceiling outcome is entirely recorded in these artifacts, reanalysis under modified scoring rules or preregistered follow-up arms (fault injection, multi-sample generation, multi-validator comparison) can be performed without re-executing the original shared-initial generation stage.

\section{Discussion}
Design implications of shared\_initial\_candidate semantics. The preregistered shared\_initial\_candidate choice isolates downstream validator effects by ensuring both conditions evaluate the same initial candidate for each pair. This isolation is scientifically valuable because any observed paired difference must derive from validator-triggered repair. It also concentrates the risk of a single-sample ceiling: if the generator sample aligns with the validator for all tasks, the design yields zero observed variance in the primary estimand, as occurred here. We emphasize this property because it determines what the paired estimator can and cannot measure: it cannot detect differences attributable to generator sampling or to alternative prompt constraints when the sampling semantics were shared.

Operational interpretation and validator scope. The deterministic validator (deterministic\_netops\_validator\_v5) reliably enforces mechanistic intent-constraint checks (parsed\_command\_count, required\_command\_count, violation\_codes, valid flag). However, artifact-grounded evidence here shows that correctness-by-structure does not imply semantic intent alignment under operator expectations. That limitation is intrinsic to validators that lack an external intent model or operator feedback; detecting semantic mismatches requires either richer intent models, human-in-the-loop signals, or tailored semantic validators. The present execution does not imply validators are unnecessary---rather, it demonstrates that in a synthetic bank and single-sample regime the validator had nothing mechanical to correct.

Prescriptions for follow-up preregistered experiments (artifact-grounded). The archived infrastructure and artifacts directly support concrete next-step preregistered arms that would produce estimable validator FP/FN and repair utility while preserving preregistration rigor:
  1) Fault-injection arm (preregistered): deterministically inject defined fault classes into a controlled fraction f of reference answers (syntactic errors, configuration-logic faults, semantic-intent swaps). This increases baseline invalidity and permits measurement of repair invocation rates and post-repair success.
  2) Multi-sample generation: set initial\_generation\_calls\_per\_task = k >= 3 independent draws per task with independent seeds recorded. Allow the guarded pipeline to accept any validator-approved draw or to perform repair. This reintroduces generator variability and creates opportunities for validator rejections and repairs.
  3) Multi-validator and multi-model comparison: add alternative validator implementations and at least one additional model family. This will quantify sensitivity to validator design and model diversity and improve external validity.

Each prescription is directly supported by the archived adapter semantics and task-generation determinism; implementing them preregistered would produce non-degenerate data suitable for estimating per-validator FP/FN, repair utility, detection latency, and predictive R\textasciicircum{}2 against an operational-impact proxy as originally planned.

Threats to validity. Internal validity: shared\_initial\_candidate intentionally isolates downstream effects but increased single-sample ceiling risk. Construct validity: validator.valid maps to mechanistic correctness but may fail to capture semantic intent misalignment. External validity: synthetic task bank, single-model (gpt-5-mini), and a single deterministic validator limit generalization to heterogeneous production networks and additional model families. Conclusion validity: the degenerate outcome yields zero variance in the primary estimand so inferential conclusions about validator benefit are not possible from this execution alone. These limitations are recorded in the execution and analysis artifacts and motivate the artifact-grounded follow-up arms described above.

\section{Limitations}
\begin{itemize}
\item Shared\_initial\_candidate semantics constrained inference: baseline and guarded reused the same initial candidate per pair; any paired difference could only arise from validator-triggered repair (not from independent generator sampling).
\item Synthetic task bank (difficulty 'very\_hard'/'extreme') is not equivalent to heterogeneous production networks (multi-vendor devices, real ticket noise); external validity is limited.
\item Single-model experiment: only gpt-5-mini was executed; no multi-model generality claims are made.
\item No repairs were applied because all initial candidates passed the deterministic validator; therefore the study could not estimate repair effectiveness or validator false-negative behavior in this execution.
\item Validator scope: deterministic\_netops\_validator\_v5 enforces mechanistic intent-constraint checks but does not guarantee detection of semantic intent misalignment (correct-by-structure but incorrect-by-intent).
\item Recorded execution/model\_configuration.json temperature = null/unset does not imply deterministic sampling; provider-side sampling behavior may vary and is documented in execution manifests.
\end{itemize}

\section{Conclusion}
We executed a fully preregistered paired confirmatory study (AC-2026-03) under shared\_initial\_candidate semantics to measure validator utility in NetOps GVR pipelines. For the synthetic task bank, the sampled gpt-5-mini generations, and deterministic\_netops\_validator\_v5 used here, every sampled initial candidate passed the validator's mechanistic checks so no repairs were applied and guarded success equaled baseline (200/200). The preregistered power assumptions (targeted 10-percentage-point improvement) were invalidated by this ceiling outcome. The result is artifact-grounded and reproducible from the archived bundle. We publish the exact execution and analysis artifacts to support replication and provide concrete, preregistered experiment prescriptions (fault injection, multi-sample generation, multi-validator comparison) that are required to produce estimable validator FP/FN behavior and repair value in operationally meaningful conditions.

\section*{Disclosure Statement}
Disclosure (concise): Declared model and provider: gpt-5-mini via provider bundle 'openai\_responses'. Orchestration/adapter: hosted\_netops\_gvr\_v1. Deterministic validator: deterministic\_netops\_validator\_v5 (validator\_version = 5.0). Preregistration: AC-2026-03 (preregistration\_sha256 = de37b485\allowbreak{}421cb089\allowbreak{}5a98df38\allowbreak{}b6b3baad\allowbreak{}1dc7c13c\allowbreak{}81065e3d\allowbreak{}8240d2a9\allowbreak{}4a0ff844). Master prompt immutable reference: provenance/master\_prompt.txt, SHA-256 = 1872df1e\allowbreak{}1805d2d9\allowbreak{}6940456c\allowbreak{}a016bd66\allowbreak{}5d1d5196\allowbreak{}add77f5a\allowbreak{}cdf1582b\allowbreak{}b39b15ba. Autonomy boundary: a human Research Director selected the broad topic family and approved the master prompt before the autonomy lock; after the lock the AI system autonomously performed literature review, preregistration, experiment design, execution, analysis, internal peer review and manuscript drafting without manual editing of technical content. Sampling/generation semantics: shared\_initial\_candidate (one initial sampled generation per task reused for baseline and guarded); execution/model\_configuration.json recorded temperature = null/unset (null/unset is not evidence of deterministic sampling and does not imply temperature = 0). Call budget and usage (preregistered): maximum\_model\_calls = 400; actual model\_calls\_used = 200 (one shared initial generation per pair). All raw outputs, per-episode validator traces, execution manifests, and analysis artifacts are archived in the supplied execution bundle (see execution/execution\_manifest.json). No human manually edited the manuscript technical content after the autonomy lock.

\begin{thebibliography}{99}
\bibitem{ref1} [1] R. Cadenas, "Convergent Correctness in Stochastic Code Generation: A Generate-Verify-Repair Architecture for Deterministic Validation of Autonomous AI Coding Agents in Regulated Environments," SSRN Electronic Journal, 2026. doi:10.2139/ssrn.6754899.
\bibitem{ref2} [2] K.-H. Tseng, N. Bogahawatta, Y. Ginige, K. Patel, K. Dakic, S. Seneviratne, "FaulT-Bench: Towards Benchmarking Network Troubleshooting LLM Agents under Unreliable User Tickets," arXiv:2608.27021, 2026.
\bibitem{ref3} [3] Y. Liu, C. Pei, L. Xu, B. Chen, M. Sun, Z. Zhang, Y. Sun, S. Zhang, K. Wang, H. Zhang, J. Li, G. Xie, X. Wen, X. Nie, M. Ma, P. Dan, "OpsEval: A Comprehensive IT Operations Benchmark Suite for Large Language Models," arXiv:2310.07637, 2023.
\bibitem{ref4} [4] R. S. Babu and A. Agrawal, "Self-Healing Agentic Orchestrators for Reliable Tool-Augmented Large Language Model Systems," arXiv:2606.01416, 2026.
\bibitem{ref5} [5] H. Zhou, C. Hu, Y. Yuan, Y. Cui, Y. Jin, C. Chen, H. Wu, D. Yuan, L. Jiang, D. Wu, X. Liu, J. Zhang, X. Wang, J. Liu, "Large Language Model (LLM) for Telecommunications: A Comprehensive Survey on Principles, Key Techniques, and Opportunities," IEEE Commun. Surveys \& Tutorials, 2024. doi:10.1109/COMST.2024.3465447.
\end{thebibliography}

\end{document}
